\documentclass[11pt,a4paper]{article}
\usepackage[T1]{fontenc}
\usepackage[utf8]{inputenc}
\usepackage{lmodern}
\usepackage[margin=27mm,headheight=14pt]{geometry}
\usepackage{amsmath,amssymb,amsthm,mathtools}
\usepackage{microtype,booktabs,array,longtable,enumitem,needspace}
\usepackage[dvipsnames]{xcolor}
\usepackage{xurl}
\usepackage[colorlinks=true,linkcolor=MidnightBlue,citecolor=MidnightBlue,urlcolor=MidnightBlue]{hyperref}
\usepackage{fancyhdr}
\pagestyle{fancy}
\fancyhf{}
\fancyhead[L]{\small Countable-action critical transseries}
\fancyhead[R]{\small Research article}
\fancyfoot[C]{\thepage}
\renewcommand{\headrulewidth}{0.3pt}
\setlength{\parskip}{4pt}
\setlength{\parindent}{1.2em}
\setlist{itemsep=3pt,topsep=5pt}
\numberwithin{equation}{section}
\newtheorem{theorem}{Theorem}[section]
\newtheorem{proposition}[theorem]{Proposition}
\newtheorem{lemma}[theorem]{Lemma}
\newtheorem{corollary}[theorem]{Corollary}
\theoremstyle{definition}
\newtheorem{definition}[theorem]{Definition}
\newtheorem{example}[theorem]{Example}
\theoremstyle{remark}
\newtheorem{remark}[theorem]{Remark}
% ed. (2026-09-29): unnumbered environment for ProveIt editorial notes; it does not shift any numbering.
\newtheorem*{ednoteinner}{Editorial note (ProveIt, 2026-09-29)}
\newenvironment{ednote}{\begin{ednoteinner}\sloppy\Urlmuskip=0mu plus 1mu\relax}{\end{ednoteinner}}
\DeclareMathOperator{\Li}{Li}
\DeclareMathOperator{\Var}{Var}
\DeclareMathOperator{\supp}{supp}
\newcommand{\N}{\mathbb N}
\newcommand{\R}{\mathbb R}
\newcommand{\C}{\mathbb C}
\newcommand{\E}{\mathbb E}
\newcommand{\Prob}{\mathbb P}
\newcommand{\dd}{\,\mathrm d}
\newcommand{\bc}{\beta_{\mathrm c}}
\newcommand{\rhoc}{\rho_{\mathrm c}}
\newcommand{\Bc}{B_{\mathrm c}}
\newcommand{\m}{\boldsymbol m}
\newcommand{\gL}{g_{\alpha,L}}
\newcommand{\Ra}{\mathcal R_\alpha}
\newcommand{\coeff}[2]{[#1]#2}
\newcommand{\repo}{\texttt{ProveIt}}
\hypersetup{pdftitle={Beyond Finite-Action Folds: Critical Hahn Transseries, Stable Sector Asymptotics, and Sharp Action Budgets},pdfauthor={Research draft prepared for Vladimir Reshetnikov},pdfsubject={Countable exponential feedback, convergent Hahn inversion, stable local limits and action cutoffs}}

\begin{document}
\begin{center}
{\small\scshape Transseries / asymptotic inversion / countable exponential feedback}\par
\vspace{10pt}
{\LARGE\bfseries Beyond Finite-Action Folds}\par
\vspace{6pt}
{\Large Critical Hahn Transseries, Stable Sector Asymptotics,\par and Sharp Action Budgets}\par
\vspace{12pt}
{\normalsize Research draft prepared for Vladimir Reshetnikov}\par
{\small 29 September 2026}\par
\end{center}
\vspace{5pt}
\begin{abstract}
Finite analytic action families generically produce square-root inverse
singularities. A countable family can instead reach a nonanalytic boundary,
where that conclusion fails. We study the explicit positive class
\[
 U_\beta(q)=\beta\sum_{j\ge1}a_jq^j e^{jU_\beta(q)},\qquad
 a_j=a j^{-1-\alpha}\quad\hbox{outside a finite set},\quad 1<\alpha<2.
\]
At its critical coupling, we construct an actually convergent Hahn inverse
with a finite positive generating grid, give a closed formula for every
coefficient, and determine its exact finite-ramification criterion.
An all-orders, uniformly additive large-sector expansion has coefficients
expressed through derivatives of a stable density. A reflection identity
identifies its terms with the Hahn coefficients and explains a precise
arithmetic cancellation rule. We then prove a joint coupling--action-cutoff
limit. The number of retained actions must exceed the stable scale
$n^{1/\alpha}$ to preserve the $n$th coefficient asymptotically; at a
comparable cutoff an explicit infinitely divisible density gives the exact
limiting loss. Finally, we compute the critical-value and curvature drift
of the finite-action folds and prove a triangular-array Gaussian limit in
the complementary hard-cutoff regime. All statements include conventional
proofs. Exact coefficient checks and reproducible numerical diagnostics
accompany the article.
\end{abstract}
\noindent\textbf{Scope and research status.}
This is a proved model-specific extension of the repository's finite-action
critical-inversion program, not a claimed solution of unrestricted
transseries realization or resurgence. Stable limits for simply generated
trees and soft/hard truncation are established subjects; those mechanisms
are credited below. Global publication priority for the combined results
has not been established. No new Lean formalization is claimed.

\clearpage
\tableofcontents
\clearpage

\section{The gap, the contribution, and the classical background}
\label{sec:intro}

\subsection{What changes when the action family is countable}
The repository contains a substantial inversion calculus and recent
research packages on support-controlled reversion, analytic moving folds,
and regularity of countable exponential feedback. The moving-fold package
constructs a local two-sheet atlas for an analytic finite-action equation;
its research questions include higher critical multiplicity, changing
actions, and the loss of a uniform action gap. The regularity package
studies $U(q)=\sum_{j\ge1}q^j e^{\lambda_jU(q)}$ and separates formal
existence, convergence, Gevrey regularity, and realization
\cite{repo-index,repo-fold,repo-regularity}.
% ed. (2026-09-29): two uncited relations to repository material, recorded on filing (batch 47).
\begin{ednote}
The finite-action result continued here is Theorem~\texttt{thm:actions} of
the moving-fold article \cite{repo-fold}. The model of
Section~\ref{sec:model}, $U=\beta\sum_j a_jq^je^{jU}$, is the slice
$c_j=\beta a_j$, $q$-exponents and slopes both equal to $j$, of the
regularity article's further-research question ``Amplitudes, general
actions, and multiple scales'' (\path{Exponential_Feedback_Regularity_Classification/article.tex},
Section~\texttt{sec:future}), which asks for $U=\sum_jc_jq^{a_j}e^{\lambda_jU}$
with general amplitudes. This article does not cite that question. For this
slice the solution converges, and the article determines its critical
boundary and exponent $1/\alpha$; the question's analogue of the exact type
identity for general $(c_j,a_j,\lambda_j)$ is not addressed.
\end{ednote}

Here the slopes are only linear, but the coefficients have a power-law
tail. The solution therefore converges near the origin, while its first
critical boundary is not analytic. The second moment diverges. The
critical exponent becomes $1/\alpha$, rather than $1/2$, and a finite
truncation of the action family restores a square root at a moving
critical value. Consequently, taking the sector index to infinity and
taking the action cutoff to infinity do not commute without a scale
condition. This is the specific gap addressed here.

We do not identify a nonanalytic boundary with an analytic critical point
of a nonintegral ``multiplicity.'' Those are different objects. Nor do we
claim to settle every higher-multiplicity or parameter-dependent-action
question in the moving-fold manuscript. Our results give a precise,
self-contained answer for the class defined in Section~\ref{sec:model}.

\subsection{What is classical, and what is supplied here}
The substitution $t=qe^U$ turns the problem into
$t=q\exp(\beta A(t))$, an instance of the classical simply generated tree
schema. Janson's survey gives critical infinite-variance local limits and
stable maximum-size behavior; in particular, Example~18.27 treats a
closely related power-law allocation setting \cite{janson}. Soft and hard
truncation of heavy tails are studied by Chakrabarty and Samorodnitsky
\cite{chakrabarty}. Lagrange inversion, singularity transfer, and saddle
methods are standard tools \cite{flajolet}. We do not relabel these general
principles as new discoveries.

The contribution developed in this article is an explicit connection
between the analytic and sector descriptions of this countable-action
model. It consists of the following linked results.
\begin{enumerate}[label=\arabic*.,leftmargin=*]
\item A convergent critical Hahn inverse, its closed multi-index formula,
its exact ramification, and a geometric truncation estimate
(Theorems~\ref{thm:hahn} and~\ref{thm:hahn-coeff}).
\item An arbitrary-order stable-density expansion and an exact
Hahn--Fourier coefficient identity, including an arithmetic cancellation
rule (Theorems~\ref{thm:allorders} and~\ref{thm:duality}).
\item A two-parameter coupling--cutoff profile and a necessary and
sufficient action-budget condition
(Theorems~\ref{thm:cutoff} and~\ref{thm:budget}).
\item Explicit finite-fold drift constants and a uniform Gaussian
coefficient law when the cutoff lies below the critical budget
(Theorems~\ref{thm:fold-drift} and~\ref{thm:hard}).
\end{enumerate}
These are intended as a reusable extension within the repository. The
proofs do not depend on the correctness of any unreviewed repository
manuscript. An independent literature and mathematical review remains
necessary before asserting publication novelty.

\section{Model, coordinates, and exact coefficients}
\label{sec:model}

Fix
\begin{equation}\label{eq:model-assumptions}
 1<\alpha<2,\qquad a>0,\qquad
 A(t)=a\Li_{1+\alpha}(t)+P(t)=\sum_{j\ge1}a_jt^j,
\end{equation}
where $P(t)=\sum_{j=1}^Jp_jt^j$, every $a_j$ is nonnegative, and $a_1>0$.
The coefficients of $P$ itself need not be nonnegative. Thus an arbitrary
finite prefix can be changed while the positive tail is held fixed. Put
\begin{equation}\label{eq:d-def}
 d_k=a\zeta(1+\alpha-k)+\sum_{j=1}^Jp_jj^k\quad(k\ge0),
 \qquad G=\Gamma(-\alpha)>0.
\end{equation}
The quantities $d_0=A(1)$ and $d_1=\sum j a_j$ are finite and positive.
For $k\ge2$, the numbers $d_k$ are \emph{analytic expansion coefficients},
not finite moments of the action distribution. In particular,
$\sum j^2a_j=\infty$, even though $d_2$ is a finite real number.

For $\beta>0$, define $U_\beta\in q\R[[q]]$ by
\begin{equation}\label{eq:feedback}
 U_\beta(q)=\beta A\bigl(qe^{U_\beta(q)}\bigr),
 \qquad u_n(\beta)=[q^n]U_\beta(q).
\end{equation}
Its coefficient at degree $n$ depends only on earlier coefficients and
$a_1,\ldots,a_n$, so the formal solution is unique. The analytic implicit
function theorem at $(q,U)=(0,0)$ also gives a convergent solution there.
All its coefficients are positive, since $a_1>0$.

Throughout, $\N=\{0,1,2,\ldots\}$. The central parameters are
\begin{equation}\label{eq:parameters}
 \bc=d_1^{-1},\quad \delta_\beta=1-\beta d_1,\quad
 \rho_\beta=e^{-\beta d_0},\quad B_\beta=\beta aG.
\end{equation}
We also write $\rhoc=\rho_{\bc}$ and $\Bc=B_{\bc}$.

\begin{proposition}[Exact parametrization and Lagrange coefficient]
\label{prop:exact}
With $t=qe^{U_\beta(q)}$ one has
\begin{equation}\label{eq:parametric}
 U_\beta=\beta A(t),\qquad q=t e^{-\beta A(t)},\qquad
 u_n(\beta)=\frac1n[t^n]e^{n\beta A(t)}\quad(n\ge1).
\end{equation}
If $N_j$ are independent Poisson random variables of means $n\beta a_j$
and $K_n=\sum_{j\ge1}jN_j$, then
\begin{equation}\label{eq:poisson}
 u_n(\beta)=\frac{\rho_\beta^{-n}}n\Prob(K_n=n),\qquad
 \E K_n=n\beta d_1.
\end{equation}
\end{proposition}
\begin{proof}
The first two identities follow directly from \eqref{eq:feedback}.
Lagrange inversion for $t=q e^{\beta A(t)}$ gives
\[
 [q^n]\beta A(t(q))
 =\frac1n[t^{n-1}]\beta A'(t)e^{n\beta A(t)}
 =\frac1n[t^n]e^{n\beta A(t)},
\]
where the last equality follows by differentiating the exponential.
The total Poisson intensity is $n\beta d_0<\infty$, so only finitely many
$N_j$ are nonzero almost surely. The probability generating function of
$K_n$ is $\exp(n\beta(A(t)-d_0))$. Comparing its coefficient at $t^n$
proves \eqref{eq:poisson}; the mean follows from $d_1<\infty$.
\end{proof}

For an integer $M\ge1$, let $A_M(t)=\sum_{j\le M}a_jt^j$ and let
$U_{\beta,M}$ and $u_{n,M}(\beta)$ denote the corresponding solution and
coefficients. A particularly useful consequence is
\begin{equation}\label{eq:conditional}
 \frac{u_{n,M}(\beta)}{u_n(\beta)}
 =\Prob\bigl(N_j=0\text{ for all }j>M\mid K_n=n\bigr).
\end{equation}
Indeed, independence contributes the factor
$\exp(-n\beta\sum_{j>M}a_j)$, and the remaining coefficient is the
truncated probability in \eqref{eq:poisson}. The ratio belongs to $[0,1]$,
is nondecreasing in $M$, and equals $1$ for $M\ge n$. This identity will
turn an analytic action-cutoff question into a conditional large-jump
question, without approximation.

\section{A convergent Hahn inverse at the critical boundary}
\label{sec:hahn}

The polylogarithm expansion \cite[Eq.~25.12.12]{dlmf}, on the branch
$|\arg x|<\pi$ and $|x|<2\pi$, gives
\begin{equation}\label{eq:polylog}
 A(e^{-x})=d_0-d_1x+aGx^\alpha
                  +\sum_{k\ge2}\frac{d_k(-x)^k}{k!}.
\end{equation}
The series on the right is convergent in that disk. Introduce
\begin{equation}\label{eq:critical-coordinate}
 x=-\log t,\qquad \varepsilon=\log(\rho_\beta/q),\qquad
 H_\beta(x)=\beta\sum_{k\ge2}\frac{d_k(-1)^k x^{k-2}}{k!}.
\end{equation}
Equations \eqref{eq:parametric} and \eqref{eq:polylog} become the exact
local normal form
\begin{equation}\label{eq:normalform}
 \boxed{\quad
 \varepsilon=\delta_\beta x+B_\beta x^\alpha+x^2H_\beta(x),
 \qquad U_\beta=\beta d_0+\varepsilon-x.\quad}
\end{equation}
For positive $q<\rhoc$ at criticality, $x$ and $\varepsilon$ are positive.
Complex branches below are continuations of this positive branch.

\begin{theorem}[Convergent critical Hahn chart]
\label{thm:hahn}
At $\beta=\bc$, set
\[
 T=(\varepsilon/\Bc)^{1/\alpha},\qquad R=T^{2-\alpha}.
\]
There is a holomorphic function $V(T,R)$ near $(0,0)$, with $V(0,0)=1$,
such that the critical inverse is exactly
\begin{equation}\label{eq:hahn-chart}
 x=T V(T,T^{2-\alpha}),\qquad
 V^\alpha+\frac{R}{\Bc}V^2H_{\bc}(TV)=1.
\end{equation}
Consequently its Hahn support in $\varepsilon$ is contained in
\begin{equation}\label{eq:hahn-grid}
 \left\{\frac{1+k+(2-\alpha)m}{\alpha}:k,m\in\N\right\}.
\end{equation}
The expansion is absolutely convergent on a sufficiently small slit
neighborhood, and its support has only finitely many elements below any
fixed bound.

More quantitatively, suppose $V$ is bounded by $K_0$ on a polydisk with
radii $T_0,R_0$. If its Taylor polynomial of total degree at most $N$ is
$V_{\le N}$, and
\[
 \theta=\max\{|T|/T_0,|T|^{2-\alpha}/R_0\}<1,
\]
then
\begin{equation}\label{eq:hahn-remainder}
 |x-TV_{\le N}(T,T^{2-\alpha})|
 \le |T|K_0\frac{(N+2)\theta^{N+1}}{(1-\theta)^2}.
\end{equation}
\end{theorem}
\begin{proof}
Substitute $x=TV$ into \eqref{eq:normalform} with $\delta_\beta=0$ and
divide by $\Bc T^\alpha$. This gives \eqref{eq:hahn-chart}. Its derivative
with respect to $V$ at $(T,R,V)=(0,0,1)$ is $\alpha\ne0$. The analytic
implicit function theorem gives the claimed $V$ and its Taylor series.
Multiplication by $T$ and then substitution $R=T^{2-\alpha}$ give
\eqref{eq:hahn-grid}. Both grid generators are positive. Thus a bound
$1+k+(2-\alpha)m\le C$ bounds $k$ and $m$ separately, proving local
finiteness even when the additive group they generate is dense.

Cauchy's estimate bounds the Taylor coefficient of $T^kR^m$ by
$K_0T_0^{-k}R_0^{-m}$. At total degree $\ell$ there are $\ell+1$ such
monomials. The omitted sum is therefore at most
\[
 |T|K_0\sum_{\ell>N}(\ell+1)\theta^\ell
 \le |T|K_0(N+2)\theta^{N+1}(1-\theta)^{-2}.
\]
This also proves absolute convergence after the substitution.
\end{proof}

\begin{remark}[Obtaining actual majorant radii]
The radii in the theorem need not remain purely existential. Choose
$0<r_0<2\pi$ and a bound $H_0\ge\max(1,\sup_{|x|\le r_0}|H_{\bc}(x)|)$.
On $|v-1|=1/4$, put $m_\alpha=\min|v^\alpha-1|>0$, using the branch
near $1$. One may take
\[
 T_0=r_0/2,\qquad
 R_0=\frac{\Bc m_\alpha}{4(5/4)^2H_0},\qquad K_0=5/4.
\]
On that $v$-circle the perturbation has modulus at most $m_\alpha/4$.
Rouch\'e's theorem gives exactly one zero, hence a simple zero depending
holomorphically on $(T,R)$, inside the circle. Slightly enlarging the
parameter radii preserves the strict inequality, so Cauchy estimates at
the displayed radii are legitimate. Evaluating $H_0$ and $m_\alpha$ with
outward rounding would produce numerical certificates; the supplied
numerical program does not perform that interval-arithmetic step.
\end{remark}

\subsection{A closed formula for every coefficient}
For a finitely supported multi-index $\m=(m_2,m_3,\ldots)$ define
\begin{equation}\label{eq:multiindex}
 m_* =\sum_{k\ge2}m_k,\qquad
 K=\sum_{k\ge2}k m_k,\qquad
 \Lambda=1+K-\alpha m_*,\qquad
 h_k=\frac{\bc d_k(-1)^k}{k!\Bc}.
\end{equation}
In particular $\Lambda\ge1+(2-\alpha)m_*>0$.

\begin{theorem}[Exact multi-index reversion]
\label{thm:hahn-coeff}
The chart in Theorem~\ref{thm:hahn} has the absolutely convergent expansion
\begin{equation}\label{eq:closed-hahn}
 x=\sum_{\m}C_{\m}\left(\prod_{k\ge2}h_k^{m_k}\right)T^\Lambda,
 \qquad
 C_{\m}=\frac{(-1)^{m_*}}{\alpha\prod_{k\ge2}m_k!}
       \frac{\Gamma((K+1)/\alpha)}{\Gamma(1+\Lambda/\alpha)}.
\end{equation}
The zero multi-index contributes $T$. Terms with equal exponent are to
be grouped; every such group is finite.
\end{theorem}
\begin{proof}
Write $s=x/T$ and $y=s^\alpha$. The equation for $y$ is
\[
 y=1-\sum_{k\ge2}z_k y^{k/\alpha},\qquad z_k=h_kT^{k-\alpha}.
\]
First retain only finitely many $z_k$ and introduce an auxiliary variable
$w$ multiplying their sum. Lagrange inversion about $y=1$, applied to
$y^{1/\alpha}$, says that the coefficient of a monomial of total degree
$m_*>0$ is
\[
 \frac{(-1)^{m_*}}{\alpha\prod m_k!}
 \left.\frac{\mathrm d^{m_*-1}}{\mathrm dy^{m_*-1}}
                 y^{(K+1)/\alpha-1}\right|_{y=1}.
\]
The derivative is
$\Gamma((K+1)/\alpha)/\Gamma((K+1)/\alpha-m_*+1)$,
which gives \eqref{eq:closed-hahn}. The formula for $m_*=0$ also yields
$1$ by $\Gamma(z+1)=z\Gamma(z)$.

A bound on $\Lambda$ bounds $m_*$ and every index $k$ appearing in $\m$,
so passing from finitely many $z_k$ to all of them is coefficientwise
finite. To justify absolute rather than just formal convergence, replace
each $h_k$ by $-|h_k|$. The same implicit-function argument applies to
the analytic majorant $\sum |h_k|T^{k-2}$. The gamma factors in
\eqref{eq:closed-hahn} are positive, so every coefficient of the resulting
majorant expansion is $|C_{\m}|\prod|h_k|^{m_k}$. Its convergence proves
absolute convergence of the original expansion and validates grouping.
\end{proof}

For example, the terms supported only on $h_2$ start as
\begin{equation}\label{eq:first-hahn}
 x=T-\frac{h_2}{\alpha}T^{3-\alpha}
       +\frac{5-\alpha}{2\alpha^2}h_2^2T^{5-2\alpha}+\cdots,
\end{equation}
while the term linear in $h_3$ is $-(h_3/\alpha)T^{4-\alpha}$.
These displayed families must be sorted together; their relative order
can depend on $\alpha$.

\begin{corollary}[Exact ramification and exponential action grid]
\label{cor:ramification}
If $\alpha=r/s$ in lowest terms, the critical inverse becomes holomorphic
after $\varepsilon=z^r$, and $r$ is the smallest possible positive integer
ramification. If $\alpha$ is irrational, no finite Puiseux cover makes
the inverse holomorphic.

If $\varepsilon=e^{-\omega X}$ with $\omega>0$ and $X\to+\infty$,
\eqref{eq:closed-hahn} is an absolutely convergent exponential Hahn
transseries. A sufficient positive action grid for $x$ has initial action
$\omega/\alpha$ and increments $\omega/\alpha$ and
$\omega(2-\alpha)/\alpha$. The analytic term $\varepsilon$ in $U$ is
included separately.
\end{corollary}
\begin{proof}
For $\alpha=r/s$, $T$ is a constant times $z^s$ and $T^{2-\alpha}$ is a
constant times $z^{2s-r}$. Both exponents are positive integers. The
leading term of $x$ is a nonzero multiple of $\varepsilon^{s/r}$, so any
finite ramification must be divisible by $r$. If $\alpha$ is irrational,
the leading exponent $1/\alpha$ cannot be made integral by multiplying
it by any positive integer. The exponential statement follows by direct
substitution in the convergent grid.
\end{proof}

\section{Coupling transition and the three coefficient regimes}
\label{sec:phase}

The fractional boundary is the transition between an inherited
power-law singularity and an interior analytic fold.

\begin{theorem}[Phase diagram]
\label{thm:phases}
For fixed $\beta>0$, the radius and leading coefficients are as follows.
If $0<\beta<\bc$, the radius is $\rho_\beta$ and
\begin{equation}\label{eq:subcritical}
 u_n(\beta)\sim\beta a\delta_\beta^{-1-\alpha}
                    \rho_\beta^{-n}n^{-1-\alpha}.
\end{equation}
At $\beta=\bc$, the radius is $\rhoc$ and
\begin{equation}\label{eq:critical-leading}
 u_n(\bc)\sim\rhoc^{-n}\Bc^{-1/\alpha}
               \frac{-1}{\Gamma(-1/\alpha)}n^{-1-1/\alpha}.
\end{equation}
If $\beta>\bc$, there is a unique $\tau\in(0,1)$ with
$\beta\tau A'(\tau)=1$. Writing
\begin{equation}\label{eq:fold-parameters}
 r_\beta=\tau e^{-\beta A(\tau)},\qquad
 v_\beta=\beta(t\partial_t)^2A(t)\big|_{t=\tau}>0,
\end{equation}
one has radius $r_\beta$ and
\begin{equation}\label{eq:supercritical}
 u_n(\beta)\sim\frac{r_\beta^{-n}}{\sqrt{2\pi v_\beta}\,n^{3/2}}.
\end{equation}
\end{theorem}
\begin{proof}
The function $\beta tA'(t)$ is strictly increasing for real $0<t<1$,
starts at $0$, and tends to $\beta d_1$. Thus $q(t)=t e^{-\beta A(t)}$
is increasing up to $t=1$ when $\beta\le\bc$, and has its unique first
maximum at $t=\tau<1$ when $\beta>\bc$.

For $\delta_\beta>0$, inversion of \eqref{eq:normalform} gives
\[
 x=\varepsilon/\delta_\beta
   -B_\beta\delta_\beta^{-1-\alpha}\varepsilon^\alpha
   +o(\varepsilon^\alpha).
\]
The first nonanalytic term of $U_\beta$ is therefore
$B_\beta\delta_\beta^{-1-\alpha}\varepsilon^\alpha$.
Singularity transfer gives its coefficient divided by $\Gamma(-\alpha)$,
which proves \eqref{eq:subcritical}. At criticality the leading
nonanalytic term is $-\Bc^{-1/\alpha}\varepsilon^{1/\alpha}$ by
Theorem~\ref{thm:hahn}, proving \eqref{eq:critical-leading}.

At the interior critical point put $t=\tau e^{-x}$ and
$\varepsilon=\log(r_\beta/q)$. Taylor expansion gives
$\varepsilon=v_\beta x^2/2+O(x^3)$ and
$U_\beta=\beta A(\tau)+\varepsilon-x$.
The singular term is $-\sqrt{2/v_\beta}\,\varepsilon^{1/2}$.
Since $\Gamma(-1/2)=-2\sqrt\pi$, this yields
\eqref{eq:supercritical}.

For completeness, these local expansions describe the actual unique
dominant singularity. Pringsheim's theorem excludes a smaller convergence
radius, since the positive real parametric branch is analytic strictly
before the stated boundary. The nonanalytic local term excludes a larger
radius. The coefficients of $t(q)$ are positive, with its
first two coefficients nonzero. Monotone convergence along the positive
radius gives a finite boundary value, either $1$ or $\tau$. At every
other point on the same circle, strictness in the triangle inequality
gives $|t(q)|$ smaller than that positive value. Consequently
$|\beta t A'(t)|<1$ there, and the implicit function theorem continues
the solution across that point. Near the positive point the local
power--fractional or square-root inverse continues in a slit sector,
by \eqref{eq:normalform} or the analytic fold expansion. Compactness of
the rest of the circle gives a common indented continuation domain.
The usual singularity-transfer contour is therefore valid
\cite[Ch.~VI]{flajolet}; no unverified competing singularity is discarded.
\end{proof}

\subsection{A uniform local coupling profile}
The same normal form gives a useful transition in the \emph{argument}
coordinate, separate from the large-sector transition below.

\begin{proposition}[Boundary-layer inverse and emerging fold]
\label{prop:local-crossover}
Suppose $\varepsilon\downarrow0$ and
\[
 \delta_\beta=\eta\Bc^{1/\alpha}\varepsilon^{(\alpha-1)/\alpha},
 \qquad \eta\text{ in a fixed compact subset of }\R.
\]
There is a unique positive solution near the critical branch, and
\begin{equation}\label{eq:local-profile}
 x=(\varepsilon/\Bc)^{1/\alpha}
   \left(v(\eta)+O\bigl(\varepsilon^{\kappa}\bigr)\right),
 \quad v(\eta)^\alpha+\eta v(\eta)=1,
 \quad \kappa=\frac{\min(\alpha-1,2-\alpha)}\alpha.
\end{equation}
The error is uniform on that compact set. As $\delta_\beta\uparrow0$
through negative values, the emerging interior fold satisfies
\begin{align}
 x_{\mathrm f}&\sim
       \left(\frac{-\delta_\beta}{\alpha\Bc}\right)^{1/(\alpha-1)},
       \label{eq:emerging-x}\\
 \log(r_\beta/\rho_\beta)&\sim
 (\alpha-1)\alpha^{-\alpha/(\alpha-1)}
 \Bc^{-1/(\alpha-1)}(-\delta_\beta)^{\alpha/(\alpha-1)}.
 \label{eq:emerging-radius}
\end{align}
\end{proposition}
\begin{proof}
Divide \eqref{eq:normalform} by $\varepsilon$, putting
$x=(\varepsilon/\Bc)^{1/\alpha}v$.
Since $B_\beta/\Bc=\beta/\bc=1-\delta_\beta$, the result is
\[
 1=\eta v+v^\alpha
      +O(\varepsilon^{(\alpha-1)/\alpha})v^\alpha
      +O(\varepsilon^{(2-\alpha)/\alpha})v^2.
\]
The positive root of $v^\alpha+\eta v=1$ is unique: for $\eta<0$ the
function first decreases below zero and then increases through $1$ only
once; for $\eta\ge0$ it increases throughout. At that root its derivative
is $1/v+(\alpha-1)v^{\alpha-1}>0$. On compact $\eta$-sets the root and
this derivative are uniformly controlled, proving \eqref{eq:local-profile}.

At the interior fold the derivative of the right side of
\eqref{eq:normalform} vanishes. Its first two terms are
$\delta_\beta+\alpha B_\beta x^{\alpha-1}$; the analytic remainder is
$O(x)=o(x^{\alpha-1})$. This proves \eqref{eq:emerging-x}. Substitution
back into the normal form gives
$\varepsilon_{\mathrm f}\sim-(\alpha-1)\Bc x_{\mathrm f}^\alpha$.
Since $\varepsilon_{\mathrm f}=\log(\rho_\beta/r_\beta)$,
\eqref{eq:emerging-radius} follows.
\end{proof}

\section{Stable densities and all-order sector asymptotics}
\label{sec:stable}

We use the principal power $(-iu)^\alpha$ for real $u$. Because
$\cos(\pi\alpha/2)<0$, the integral
\begin{equation}\label{eq:stable-density}
 g_\alpha(z)=\frac1{2\pi}\int_{\R}
          \exp\bigl((-iu)^\alpha-iuz\bigr)\dd u
\end{equation}
converges absolutely, together with every derivative in $z$. This is the
density of a centered, spectrally positive stable law in a normalization
chosen to match the feedback equation.

\begin{lemma}[Normalization, derivatives, and positivity]
\label{lem:stable}
The exponent in \eqref{eq:stable-density} has the L\'evy representation
\begin{equation}\label{eq:levy}
 (-iu)^\alpha=\frac1G\int_0^\infty
       (e^{iux}-1-iux)x^{-1-\alpha}\dd x.
\end{equation}
The function $g_\alpha$ is a smooth probability density, positive for
every real argument. At the origin,
\begin{equation}\label{eq:stable-derivatives}
 g_\alpha^{(k)}(0)=\frac{\Gamma((k+1)/\alpha)}{\pi\alpha}
                   \sin\frac{\pi(k+1)}\alpha,
 \qquad
 g_\alpha(0)=-\frac1{\Gamma(-1/\alpha)}>0.
\end{equation}
\end{lemma}
\begin{proof}
For $\Re s>0$, two integrations by parts give
\[
 \int_0^\infty(e^{-sx}-1+sx)x^{-1-\alpha}\dd x
 =\frac{\Gamma(2-\alpha)}{\alpha(\alpha-1)}s^\alpha
 =\Gamma(-\alpha)s^\alpha.
\]
Pass to $s=-iu$ by continuity; integrability near zero follows from the
quadratic cancellation, and at infinity from $\alpha>1$. This proves
\eqref{eq:levy}. The centered compensated Poisson construction on
$(0,1]$, followed by the finite-intensity jumps on $(1,\infty)$, realizes
this characteristic function. Fourier integrability then gives the
smooth density in \eqref{eq:stable-density}.

For the derivative formula, restrict the Fourier integral to $u>0$ and
take twice its real part. With $c=-e^{-i\pi\alpha/2}$, which has positive
real part, substitute $w=u^\alpha$ to obtain
\[
 g_\alpha^{(k)}(0)=\frac{\Gamma((k+1)/\alpha)}{\pi\alpha}
       \Re\bigl((-i)^k c^{-(k+1)/\alpha}\bigr).
\]
The real part is $\sin(\pi(k+1)/\alpha)$, giving the first identity.
Gamma reflection yields the second.

Here is also a direct positivity argument useful later. For any $L>0$
let $X_L$ have the compensated L\'evy measure
$G^{-1}x^{-1-\alpha}\mathbf1_{(0,L]}(x)\dd x$.
Its density exists by the Fourier estimates in
Lemma~\ref{lem:truncated-density} below. Split it at $\epsilon<L$:
\[
 X_L=X_\epsilon+S_{\epsilon,L}
                  -\frac1G\int_\epsilon^L x^{-\alpha}\dd x,
\]
where $S_{\epsilon,L}$ is an independent compound Poisson sum.
The variance of $X_\epsilon$ tends to zero, while the displayed
compensation tends to positive infinity as $\epsilon\downarrow0$.
Conditional on a sufficiently large finite number of jumps,
$S_{\epsilon,L}$ has a density positive throughout the interval between
that number times $\epsilon$ and that number times $L$. These intervals
overlap for large numbers of jumps. For any prescribed real $z$, choose
$\epsilon$ small enough that $z$ plus the compensation lies strictly
inside such an interval. Convolution with the positive-probability event
that $X_\epsilon$ is close to zero then gives a strictly positive density
at $z$. Thus $g_{\alpha,L}(z)>0$. Adding the jumps above $L$ to $X_L$,
and using the positive probability of no such jumps, proves
$g_\alpha(z)>0$. This argument does not assume positivity merely from
Fourier integrability.
\end{proof}

For a coupling $\beta$, which may now depend on $n$, put
\begin{equation}\label{eq:bn-lambda}
 b_n=(n\beta aG)^{1/\alpha},\qquad
 \lambda_n=\frac{n(1-\beta d_1)}{b_n},\qquad
 c_{k,n}=\frac{n\beta d_k}{k!b_n^k}\quad(k\ge2).
\end{equation}
The natural critical coupling window is $\lambda_n=O(1)$, equivalently
$\beta-\bc=O(n^{-(\alpha-1)/\alpha})$.

\begin{theorem}[Uniform all-order stable expansion]
\label{thm:allorders}
Fix a compact interval $I\subset(0,\infty)$ for $\beta$ and $R>0$. For a
finite multi-index $\m$ as in \eqref{eq:multiindex}, put
\begin{equation}\label{eq:weight}
 w(\m)=\frac{K-\alpha m_*}{\alpha}
       =\sum_{k\ge2}\frac{k-\alpha}{\alpha}m_k.
\end{equation}
Uniformly for $\beta\in I$, as $n\to\infty$,
\begin{equation}\label{eq:allorders}
 \boxed{\quad
 n b_n\rho_\beta^n u_n(\beta)
 =\sum_{w(\m)<R}
      \frac{\prod_{k\ge2}c_{k,n}^{m_k}}{\prod_{k\ge2}m_k!}
      (-1)^K g_\alpha^{(K)}(\lambda_n)
      +O_{I,R}(n^{-R}).\quad}
\end{equation}
The sum is finite. The error is an \emph{absolute, additive} error in the
displayed normalization; it is not a relative estimate uniformly far
from the critical window. In particular, if $\lambda_n\to\lambda\in\R$,
\begin{equation}\label{eq:leading-window}
 u_n(\beta)\sim\frac{\rho_\beta^{-n}}{n b_n}g_\alpha(\lambda).
\end{equation}
\end{theorem}
\begin{proof}
We provide the Fourier estimates explicitly. By
Proposition~\ref{prop:exact}, the left side of \eqref{eq:allorders} is
$b_n\Prob(K_n=n)$. Fourier inversion on the lattice gives
\begin{align}
 b_n\Prob(K_n=n)
 &=\frac1{2\pi}\int_{-\pi b_n}^{\pi b_n}
         \chi_n(u)e^{-i\lambda_nu}\dd u,\label{eq:lattice-fourier}\\
 \chi_n(u)
 &=\exp\left(n\beta\left[A(e^{iu/b_n})-d_0
                                    -id_1u/b_n\right]\right).\nonumber
\end{align}
For $|u|/b_n<2\pi$, \eqref{eq:polylog} implies the exact local expansion
\begin{equation}\label{eq:characteristic-expansion}
 \chi_n(u)=\exp\left((-iu)^\alpha
                     +\sum_{k\ge2}c_{k,n}(iu)^k\right).
\end{equation}

The needed bounds on the rest of the Fourier contour follow from the
positive tail, not from a formal Taylor series. There exist $c>0$ and
$\theta_0>0$, independent of $\beta\in I$, such that
\begin{equation}\label{eq:fourier-bound}
 \sum_{j\ge1}a_j(1-\cos(j\theta))\ge
 \begin{cases}
 c|\theta|^\alpha,&0<|\theta|\le\theta_0,\\
 c,&\theta_0\le|\theta|\le\pi.
 \end{cases}
\end{equation}
For the first bound, sum only over a block of integers with
$1\le j|\theta|\le2$; they are in the unchanged tail when $\theta_0$ is
small. Each cosine defect is bounded below, and the sum of $j^{-1-\alpha}$
over that block is comparable to $|\theta|^\alpha$. For the second bound,
the single coefficient $a_1>0$ suffices on the compact interval shown.
Thus $|\chi_n(u)|\le e^{-c'|u|^\alpha}$ for
$|u|\le\theta_0b_n$, and $|\chi_n(u)|\le e^{-c'n}$ beyond that interval.

Choose a central cutoff $L_n=(C\log n)^{1/\alpha}$, with $C$ large enough
in terms of $R$. The contribution from $|u|>L_n$ to
\eqref{eq:lattice-fourier} is $O(n^{-R})$, and so are the tails of all
finitely many stable-density derivative integrals used below.
On $|u|\le L_n$, the analytic perturbation in
\eqref{eq:characteristic-expansion} is
$O(n^{-(2-\alpha)/\alpha}|u|^2)$ and tends to zero uniformly.
The analytic series can be cut after an index $k_0$ with
$(k_0-\alpha)/\alpha>R$. Cauchy bounds for its coefficients show that
the omitted part is bounded by a constant times
$n^{-R}(1+|u|^{k_0+1})$ on this central interval, after increasing $k_0$
if necessary.

Expand the exponential of the remaining finite polynomial. Each
multi-index contributes exactly
$\prod c_{k,n}^{m_k}(iu)^K/\prod m_k!$, with size
$O(n^{-w(\m)}|u|^K)$. Because every basic weight
$(k-\alpha)/\alpha$ is positive, only finitely many have weight below
$R$. Taylor's remainder and the finite analytic coefficient bounds leave
an error of the form
\[
 C_R n^{-R}(1+|u|^{D_R})e^{-c''|u|^\alpha}
\]
after multiplication by the stable factor, for a finite integer $D_R$.
This bound is integrable with a constant independent of $n$ and $\beta$.
It proves the asserted remainder without introducing unwanted powers of
$\log n$.

Finally,
\[
 \frac1{2\pi}\int_\R(iu)^K
         e^{(-iu)^\alpha-i\lambda_nu}\dd u
       =(-1)^K g_\alpha^{(K)}(\lambda_n).
\]
In all estimates the oscillatory factor $e^{-i\lambda_nu}$ has modulus
one, so they are uniform even when $\lambda_n$ is unbounded. This proves
the additive statement. For bounded $\lambda_n$, positivity and
continuity of $g_\alpha$ give a positive lower bound on compact intervals,
which justifies the relative conclusion \eqref{eq:leading-window}.
\end{proof}

\begin{remark}[Why the qualification on the error matters]
For fixed subcritical or supercritical $\beta$, $|\lambda_n|$ diverges.
The leading density can then be smaller than the additive remainder in
\eqref{eq:allorders}. One cannot obtain a uniform relative theorem in
those regimes merely by dividing that equation by $g_\alpha(\lambda_n)$.
The fixed-coupling results in Theorem~\ref{thm:phases} have a separate
singularity proof. A quantitative matching theorem on an expanding
$\lambda$-window is left as a research problem.
\end{remark}

\section{An exact Hahn--Fourier identity and the cancellation sieve}
\label{sec:duality}

The local inverse and the sector expansion are not two unrelated
approximations. Their coefficients match term by term.

\begin{theorem}[Coefficient duality]
\label{thm:duality}
At criticality, the multi-index $\m$ contributes to $x$ the term
\[
 C_{\m}\left(\prod h_k^{m_k}\right)
                    \Bc^{-\Lambda/\alpha}\varepsilon^{\Lambda/\alpha}.
\]
Its nonanalytic contribution to $u_n$ has multiplier
\begin{equation}\label{eq:duality-left}
 -\frac{C_{\m}\prod h_k^{m_k}\,\Bc^{-\Lambda/\alpha}}
          {\Gamma(-\Lambda/\alpha)}
\end{equation}
in front of $\rhoc^{-n}n^{-1-\Lambda/\alpha}$. This multiplier equals
\begin{equation}\label{eq:duality-right}
 \frac{(-1)^K g_\alpha^{(K)}(0)}{\prod m_k!}
       \bc^{m_*}\Bc^{-(K+1)/\alpha}
       \prod_{k\ge2}\left(\frac{d_k}{k!}\right)^{m_k},
\end{equation}
which is exactly the corresponding term of
Theorem~\ref{thm:allorders} after division by $n b_n$.

The universal gamma factor vanishes precisely when
$(K+1)/\alpha$ is an integer. If $\alpha=r/s$ in lowest terms this means
$r\mid K+1$. Equivalently, that Hahn term is an integer power of
$\varepsilon$ and is locally analytic. For irrational $\alpha$ this
particular cancellation never occurs. Other cancellations caused by
zero $d_k$ or by grouping equal exponents are not excluded.
\end{theorem}
\begin{proof}
In the logarithmic coordinate $q=\rhoc e^{-\varepsilon}$, Cauchy's
coefficient differential is exactly
$q^{-n}\dd q/q=-\rhoc^{-n}e^{n\varepsilon}\dd\varepsilon$.
The local Hankel integral for a power $\varepsilon^z$ therefore has
factor $n^{-z-1}/\Gamma(-z)$. There is no additional expansion of a
Jacobian. This logarithmic transfer can also be checked against the
polylogarithm identity \eqref{eq:polylog}, whose coefficients are exact
powers of $n$. The minus sign in \eqref{eq:duality-left} comes from
$U=\bc d_0+\varepsilon-x$.

For the algebraic identity, write $z=\Lambda/\alpha=(K+1)/\alpha-m_*$.
Gamma reflection gives
\[
 \frac1{\Gamma(1+z)\Gamma(-z)}=-\frac{\sin\pi z}{\pi},
 \qquad
 \sin\pi z=(-1)^{m_*}\sin\frac{\pi(K+1)}\alpha.
\]
Insert \eqref{eq:closed-hahn} and
$h_k=\bc d_k(-1)^k/(k!\Bc)$ into \eqref{eq:duality-left}.
The resulting sign is $(-1)^K$, the remaining gamma--sine factor is
$g_\alpha^{(K)}(0)$ by \eqref{eq:stable-derivatives}, and the power of
$\Bc$ is
$-m_*-\Lambda/\alpha=-(K+1)/\alpha$. This is
\eqref{eq:duality-right}.

Since $\Lambda>0$, the reciprocal gamma factor in
\eqref{eq:duality-left} vanishes exactly when $\Lambda/\alpha$ is a
positive integer. This is equivalent to $(K+1)/\alpha\in\N$.
For $\alpha=r/s$ the condition is $r\mid K+1$. The leading monomial
coefficients are understood with continuous reciprocal gamma at its
zeros, so the identity includes the vanishing cases.
\end{proof}

\begin{example}[The first correction disappears at $\alpha=3/2$]
\label{ex:threehalves}
For $\alpha=3/2$, $g_\alpha''(0)=0$, since $(2+1)/\alpha=2$.
Consequently the nominal relative $n^{-1/3}$ correction
$c_{2,n}g_\alpha''(0)$ is absent at critical coupling, regardless of the
finite prefix. If $d_2\ne0$, the first possible nonzero correction is
\begin{equation}\label{eq:threehalf-correction}
 n b_n\rhoc^n u_n(\bc)
  =g_{3/2}(0)+\frac{c_{2,n}^2}{2}g_{3/2}^{(4)}(0)+O(n^{-1}).
\end{equation}
It has order $n^{-2/3}$. The corresponding first Hahn correction is
$-(h_2/\alpha)T^{3-\alpha}=-(h_2/\alpha)T^{3/2}$, an integer power of
$\varepsilon$; hence it contributes no singular large-sector term.
This explains the cancellation structurally, rather than discovering it
only by a vanishing numerical constant.
\end{example}

\begin{remark}[Two different kinds of convergence]
The critical Hahn expansion in $\varepsilon$ is convergent. The expansion
in decreasing powers of $n$ in Theorem~\ref{thm:allorders} is asserted to
all finite orders, not to converge when infinitely many orders are summed.
The derivative factors $g_\alpha^{(K)}$ can grow rapidly with $K$.
The identity between their individual coefficients does not identify the
two notions of convergence and does not imply Borel summability of the
large-sector expansion.
\end{remark}

\section{A sharp joint coupling--action-cutoff profile}
\label{sec:cutoff}

For $L>0$, define
\begin{align}
 \Psi_{\alpha,L}(u)
 &=\frac1G\int_0^L(e^{iux}-1-iux)x^{-1-\alpha}\dd x,
       \label{eq:truncated-exponent}\\
 g_{\alpha,L}(z)
 &=\frac1{2\pi}\int_\R e^{\Psi_{\alpha,L}(u)-iuz}\dd u,
       \label{eq:truncated-density}\\
 \nu_L&=\frac1{\alpha G L^\alpha},\qquad
 h_L=\frac{L^{1-\alpha}}{(\alpha-1)G}.
       \label{eq:tail-constants}
\end{align}
Here $\nu_L$ is the intensity of omitted jumps above $L$, and $h_L$ is
their mean in stable-scale units. Both constants are necessary: deleting
large jumps changes the mass and the centering.

\begin{lemma}[Truncated density estimates]
\label{lem:truncated-density}
The function $g_{\alpha,L}$ is a smooth, strictly positive probability
density. For $L$ in a compact subinterval of $(0,\infty)$, its defining
Fourier integral and all fixed derivatives admit common integrable
majorants. As $L\to\infty$, $g_{\alpha,L}\to g_\alpha$ locally uniformly.
For $0<L\le1$ there is a constant depending only on $\alpha$ such that
\begin{equation}\label{eq:small-L-density}
 \sup_z g_{\alpha,L}(z)\le C_\alpha L^{-(2-\alpha)/2}.
\end{equation}
\end{lemma}
\begin{proof}
The compensated Poisson construction and the positivity argument in
Lemma~\ref{lem:stable} apply with the L\'evy measure restricted to $(0,L]$.
To justify the density used there, observe that
\[
 -\Re\Psi_{\alpha,L}(u)
   =\frac1G\int_0^L(1-\cos ux)x^{-1-\alpha}\dd x
 \ge c_\alpha
 \begin{cases}
 u^2L^{2-\alpha},& |u|L\le1,\\
 |u|^\alpha,& |u|L\ge1.
 \end{cases}
\]
The first inequality uses $1-\cos v\ge c v^2$ on $|v|\le1$.
For the second, restrict to $0<x<1/|u|$ and use the first estimate there.
These bounds prove integrability with any polynomial factor, hence
smoothness and the uniform majorants. They also give
\[
 \int_\R |e^{\Psi_{\alpha,L}(u)}|\dd u
 \le C L^{-(2-\alpha)/2}+C,
\]
which proves \eqref{eq:small-L-density}. As $L\to\infty$, the exponent
converges to \eqref{eq:levy}; for $L\ge1$ the same lower bounds give a
common integrable majorant. Dominated convergence proves the local
uniform convergence of the densities.
\end{proof}

\Needspace{15\baselineskip}
\begin{theorem}[Joint critical cutoff law]
\label{thm:cutoff}
Let $\beta=\beta_n$ stay in a compact subinterval of $(0,\infty)$, and
use $b_n,\lambda_n$ from \eqref{eq:bn-lambda}. Suppose
\[
 \lambda_n\longrightarrow\lambda\in\R,\qquad
 \frac{M_n}{b_n}\longrightarrow L\in(0,\infty).
\]
Then
\begin{equation}\label{eq:cutoff-profile}
 \boxed{\quad
 \frac{u_{n,M_n}(\beta_n)}{u_n(\beta_n)}
 \longrightarrow
 \Ra(L,\lambda):=
 e^{-\nu_L}\frac{g_{\alpha,L}(\lambda+h_L)}{g_\alpha(\lambda)}.
 \quad}
\end{equation}
The convergence is uniform for $M_n/b_n$ and $\lambda_n$ in fixed compact
subintervals of $(0,\infty)$ and $\R$, respectively, in the equivalent
form with $\Ra(M_n/b_n,\lambda_n)$ on the right and an $o(1)$ error.
The profile is independent of $a$ and of the finite prefix $P$ after
these quantities have been incorporated into $\bc$ and $b_n$.
\end{theorem}
\begin{proof}
Let $K_{n,M}=\sum_{j\le M}jN_j$ and
$\mu_{n,M}=n\beta\sum_{j\le M}j a_j$. The exact conditional identity
can be written
\begin{equation}\label{eq:cutoff-ratio-exact}
 \frac{u_{n,M}(\beta)}{u_n(\beta)}
 =\exp\left(-n\beta\sum_{j>M}a_j\right)
       \frac{\Prob(K_{n,M}=n)}{\Prob(K_n=n)}.
\end{equation}
Because $M\sim Lb_n\to\infty$, the omitted weights are in the unchanged
tail. Integral comparison gives, uniformly on compact $L$-intervals,
\begin{align}
 n\beta\sum_{j>M}a_j&\longrightarrow\nu_L,\label{eq:cutoff-mass}\\
 \frac{n\beta\sum_{j>M}j a_j}{b_n}&\longrightarrow h_L.
       \label{eq:cutoff-mean}
\end{align}
The target centered lattice coordinate is consequently
$(n-\mu_{n,M})/b_n=\lambda_n+h_L+o(1)$.

The centered characteristic exponent of $K_{n,M}/b_n$ is
\[
 n\beta\sum_{j\le M}a_j
        (e^{iu j/b_n}-1-iu j/b_n).
\]
Riemann summation, using $n\beta a/b_n^\alpha=1/G$, shows its convergence
to $\Psi_{\alpha,L}(u)$, locally uniformly in $u$ and $L$.
The integrand at the lower endpoint is integrable after its quadratic
cancellation. Explicitly, the contribution of $j\le\epsilon b_n$ is
bounded by $C u^2\epsilon^{2-\alpha}$, uniformly in $n$; on the remainder
ordinary Riemann convergence applies. The finitely changed prefix
contributes only $O(n/b_n^2)=o(1)$ to the centered exponent.

To upgrade weak convergence to the required lattice local limit, one
must control the Fourier tails. For $M$ sufficiently large, unchanged
tail weights give
\begin{equation}\label{eq:trunc-fourier-bound}
 \sum_{j\le M}a_j(1-\cos j\theta)\ge c
 \begin{cases}
 \theta^2M^{2-\alpha},&|\theta|M\le1,\\
 |\theta|^\alpha,&1/M\le|\theta|\le\theta_0.
 \end{cases}
\end{equation}
For the first bound use a terminal block of integers below $M$ and the
quadratic cosine estimate. For the second, use integers up to
$1/|\theta|$; after making $\theta_0$ small the unchanged tail supplies
a fixed fraction of their second-moment sum. The coefficient $a_1$
again supplies a uniform gap for $\theta_0\le|\theta|\le\pi$.
Upon putting $u=b_n\theta$ and restricting $M/b_n$ to a compact
positive interval, these inequalities give a common integrable
majorant for the rescaled characteristic functions, with an
exponentially small contribution from the macroscopic arc.
Fourier inversion and dominated convergence therefore yield
\[
 b_n\Prob(K_{n,M}=n)
   =g_{\alpha,M/b_n}\left(\lambda_n+
                   \frac{n\beta\sum_{j>M}j a_j}{b_n}\right)+o(1)
   \longrightarrow g_{\alpha,L}(\lambda+h_L).
\]
All bounds are uniform on the compact parameter sets specified in the
theorem. The denominator is
$b_n\Prob(K_n=n)=g_\alpha(\lambda_n)+o(1)$ by
Theorem~\ref{thm:allorders}, and it stays bounded away from zero there.
Insert these estimates and \eqref{eq:cutoff-mass} into
\eqref{eq:cutoff-ratio-exact}. This proves \eqref{eq:cutoff-profile} and
its uniform version. The disappearance of the finite prefix from the
centered Riemann limit proves the universality assertion.
\end{proof}

\subsection{Why the profile is a genuine conditional probability}
The L\'evy measure above $L$ has finite mass $\nu_L$ and first moment
$h_L$. Decomposing the stable variable into jumps below and above $L$
gives the nonnegative density identity
\begin{equation}\label{eq:poisson-density-decomposition}
 g_\alpha(z)=e^{-\nu_L}\sum_{m=0}^\infty\frac1{m!}
   \int_{(L,\infty)^m}
      g_{\alpha,L}\left(z+h_L-\sum_{i=1}^m x_i\right)
      \prod_{i=1}^m\frac{\dd x_i}{Gx_i^{1+\alpha}}.
\end{equation}
For $m=0$, the integral is understood to be $g_{\alpha,L}(z+h_L)$.
Thus \eqref{eq:cutoff-profile} is exactly the fraction of the stable
density arising from the event of no jump above $L$, when conditioned on
the final centered value. In particular
\begin{equation}\label{eq:profile-range}
 0<\Ra(L,\lambda)<1\qquad(L>0,\ \lambda\in\R).
\end{equation}
The first inequality follows from density positivity. The second is
strict because the $m=1$ term in \eqref{eq:poisson-density-decomposition}
is positive. Ignoring either $e^{-\nu_L}$ or the shift $h_L$ would
produce the wrong profile.

\begin{theorem}[Necessary and sufficient critical action budget]
\label{thm:budget}
Fix $\beta=\bc$ and put $b_n=(n\Bc)^{1/\alpha}$. For every sequence of
positive integer cutoffs $M_n$,
\begin{equation}\label{eq:budget-iff}
 \frac{u_{n,M_n}(\bc)}{u_n(\bc)}\longrightarrow1
 \quad\Longleftrightarrow\quad \frac{M_n}{b_n}\longrightarrow\infty.
\end{equation}
If $M_n/b_n\to0$, the coefficient ratio tends to zero. If
$M_n/b_n\to L\in(0,\infty)$, it tends to $\Ra(L,0)\in(0,1)$.
The function $L\mapsto\Ra(L,\lambda)$ is nondecreasing and satisfies
\begin{equation}\label{eq:profile-endpoints}
 \lim_{L\downarrow0}\Ra(L,\lambda)=0,\qquad
 \lim_{L\to\infty}\Ra(L,\lambda)=1
\end{equation}
for each fixed $\lambda$.
\end{theorem}
\begin{proof}
Monotonicity follows by applying \eqref{eq:conditional} at two cutoffs
$\lfloor L_1b_n\rfloor\le\lfloor L_2b_n\rfloor$ and passing to the
limit. For $\lambda\ne0$, choose the unique nearby coupling satisfying
$\lambda_n=\lambda$; the defining expression in
\eqref{eq:bn-lambda} is strictly decreasing in $\beta$, so this choice
exists for all large $n$.

As $L\to\infty$, $h_L,\nu_L\to0$ and
$g_{\alpha,L}\to g_\alpha$ locally uniformly by
Lemma~\ref{lem:truncated-density}. This proves the second limit in
\eqref{eq:profile-endpoints}. For $L\downarrow0$, the same lemma gives
\[
 \Ra(L,\lambda)\le
  \frac{C_\alpha}{g_\alpha(\lambda)}
  L^{-(2-\alpha)/2}\exp\left(-\frac1{\alpha G L^\alpha}\right)
  \longrightarrow0.
\]

If $M_n/b_n\to\infty$, compare the coefficient ratio below with the
ratio at $\lfloor Lb_n\rfloor$ for each fixed $L$, and above with $1$.
Theorem~\ref{thm:cutoff}, followed by $L\to\infty$, proves sufficiency.
If $M_n/b_n\to0$, compare above with the ratio at
$\lfloor Lb_n\rfloor$ and then let $L\downarrow0$.
Finally, if $M_n/b_n$ does not tend to infinity, a subsequence is bounded
above by some finite $L$. On that subsequence the upper limit of the
coefficient ratio is at most $\Ra(L,0)<1$ by
\eqref{eq:profile-range}. Thus a ratio tending to $1$ forces
$M_n/b_n\to\infty$. The comparable-cutoff case is
Theorem~\ref{thm:cutoff} itself.
\end{proof}

\begin{remark}[What the budget does and does not mean]
Every coefficient of degree $n$ is reproduced \emph{exactly} by retaining
$n$ actions. The theorem identifies the much smaller scale needed for
asymptotic relative accuracy, but requires a factor tending to infinity
beyond $n^{1/\alpha}$ for the ratio to tend to $1$. A fixed large multiple
can give a very small limiting error, but not zero limiting error.
This is a theorem about growing sector index, not about truncating a
convergent sum at a fixed numerical value of $q$.
\end{remark}

\section{How finite-action folds approach the fractional boundary}
\label{sec:finite-folds}

At critical coupling every finite truncation has finite variance and an
ordinary fold. The location and curvature of that fold contain the same
scale that appears in the action-budget theorem.

Let $M\ge J$, and let $\tau_M>1$ be the unique positive solution of
\begin{equation}\label{eq:finite-critical}
 \bc\sum_{j\le M}j a_j\tau_M^j=1.
\end{equation}
It is larger than $1$ because the left side at $1$ is less than
$\bc d_1=1$. Define
\begin{equation}\label{eq:finite-r-v}
 r_M=\tau_Me^{-\bc A_M(\tau_M)},\qquad
 v_M=\bc\sum_{j\le M}j^2a_j\tau_M^j.
\end{equation}
In particular, $r_M>e^{-\bc A_M(1)}>\rhoc$.

There is a unique positive number $v_\alpha$ satisfying
\begin{equation}\label{eq:v-alpha}
 \int_0^1(e^{v_\alpha u}-1)u^{-\alpha}\dd u
                    =\frac1{\alpha-1}.
\end{equation}
Existence and uniqueness follow because the left side is continuous,
strictly increasing from $0$ to infinity. Introduce
\begin{align}
 J_\alpha(v)&=\frac1\alpha+\frac{v}{\alpha-1}
       -\int_0^1(e^{vu}-1-vu)u^{-1-\alpha}\dd u,
       \label{eq:J-alpha}\\
 D_\alpha&=\bc a J_\alpha(v_\alpha),\qquad
 V_\alpha=\bc a\int_0^1u^{1-\alpha}e^{v_\alpha u}\dd u.
       \label{eq:D-V}
\end{align}
The integrals converge at zero. Moreover,
$J_\alpha'(v_\alpha)=0$ and $J_\alpha''(v)<0$.
Since $J_\alpha'(0)=1/(\alpha-1)>0$, one has
$J_\alpha(v_\alpha)>J_\alpha(0)=1/\alpha$, so $D_\alpha>0$.

\begin{theorem}[Critical-value and curvature drift]
\label{thm:fold-drift}
As $M\to\infty$,
\begin{equation}\label{eq:fold-drift}
 M\log\tau_M\longrightarrow v_\alpha,\qquad
 \log(r_M/\rhoc)\sim D_\alpha M^{-\alpha},\qquad
 v_M\sim V_\alpha M^{2-\alpha}.
\end{equation}
The number $v_\alpha$ depends only on $\alpha$, and the dependence of
$D_\alpha,V_\alpha$ on the finite prefix is only through $\bc a$.
\end{theorem}
\begin{proof}
Write $h_M=\log\tau_M$. Subtract
$\bc\sum_{j\ge1}j a_j=1$ from \eqref{eq:finite-critical} to get
\begin{equation}\label{eq:h-root}
 \sum_{j\le M}j a_j(e^{h_Mj}-1)=\sum_{j>M}j a_j.
\end{equation}
For any fixed $v\ge0$, multiply the difference between the two sides,
with $h_M$ replaced by $v/M$, by $M^{\alpha-1}$. Tail summation and a
Riemann sum give
\[
 a\left(\int_0^1(e^{vu}-1)u^{-\alpha}\dd u
                         -\frac1{\alpha-1}\right).
\]
The finite prefix contributes $O(M^{\alpha-2})=o(1)$. The convergence is
uniform for $v$ on compact intervals. Strict monotonicity in $h$ brackets
the root between $(v_\alpha-\epsilon)/M$ and
$(v_\alpha+\epsilon)/M$ for every fixed $\epsilon>0$, proving the first
limit in \eqref{eq:fold-drift}.

An exact rearrangement, important also for stable numerical evaluation,
is
\begin{equation}\label{eq:radius-rearrangement}
 \log(r_M/\rhoc)=\bc\left(
       h_M\sum_{j>M}j a_j+\sum_{j>M}a_j
       -\sum_{j\le M}a_j(e^{h_Mj}-1-h_Mj)\right).
\end{equation}
Indeed, begin with
$h_M+\bc(d_0-A_M(e^{h_M}))$ and use $1=\bc d_1$.
After multiplying by $M^\alpha$, the first two terms in the brackets
tend to $a v_\alpha/(\alpha-1)$ and $a/\alpha$.
The last term tends to
$a\int_0^1(e^{v_\alpha u}-1-v_\alpha u)u^{-1-\alpha}\dd u$.
Its finite-prefix contribution is $O(M^{\alpha-2})$, again negligible.
This proves the radius drift with constant $D_\alpha$.

Finally, Riemann summation in \eqref{eq:finite-r-v} gives
\[
 M^{\alpha-2}v_M
   \longrightarrow\bc a\int_0^1u^{1-\alpha}e^{v_\alpha u}\dd u
   =V_\alpha.
\]
The integrability at zero permits the same small-initial-block argument
as in the proof of Theorem~\ref{thm:cutoff}; finite-prefix terms vanish.
\end{proof}

For fixed $M$, the analytic fold law is
\begin{equation}\label{eq:fixed-M}
 u_{n,M}(\bc)\sim\frac{r_M^{-n}}{\sqrt{2\pi v_M}\,n^{3/2}}.
\end{equation}
The next theorem allows $M$ to increase with $n$, and identifies the
range in which that formula remains valid with its exact moving
parameters.

\begin{theorem}[Uniform fold law in the hard-cutoff regime]
\label{thm:hard}
If $M=M_n\to\infty$ and $n/M_n^\alpha\to\infty$, then
\begin{equation}\label{eq:hard}
 u_{n,M_n}(\bc)=
 \frac{r_{M_n}^{-n}}{\sqrt{2\pi v_{M_n}}\,n^{3/2}}(1+o(1)).
\end{equation}
Thus the finite-fold coefficient law is uniform whenever
$M_n=o(n^{1/\alpha})$. At the critical scale $M_n\asymp n^{1/\alpha}$,
the non-Gaussian profile of Theorem~\ref{thm:cutoff} is required instead.
\end{theorem}
\begin{proof}
Exponentially tilt the coefficient identity at $t=\tau_M$. Let
$\widetilde N_j$ be independent Poisson variables of means
$n\bc a_j\tau_M^j$, $j\le M$, and set
$\widetilde K=\sum_{j\le M}j\widetilde N_j$. Then
\begin{equation}\label{eq:tilted-identity}
 u_{n,M}(\bc)=\frac{r_M^{-n}}n\Prob(\widetilde K=n),\qquad
 \E\widetilde K=n,\qquad \Var(\widetilde K)=n v_M.
\end{equation}
Put $\sigma_M=\sqrt{n v_M}$ and $H_n=n/M^\alpha$. By
Theorem~\ref{thm:fold-drift},
\[
 \frac{M}{\sigma_M}\asymp H_n^{-1/2}\longrightarrow0,
 \qquad \tau_M^M\longrightarrow e^{v_\alpha}.
\]
The third absolute cumulant bound is
\[
 \frac{n\bc\sum_{j\le M}j^3a_j\tau_M^j}{\sigma_M^3}
 =O(H_n^{-1/2})=o(1).
\]
Consequently the centered characteristic function, at argument
$u/\sigma_M$, converges to $e^{-u^2/2}$ on every bounded $u$-interval.

For the local, rather than only weak, central limit, split the Fourier
integral into three regions. When $|\theta|\le1/M$, the quadratic cosine
bound gives
\[
 \left|\E e^{i\theta(\widetilde K-n)}\right|
       \le e^{-c n v_M\theta^2}.
\]
After $u=\sigma_M\theta$, this supplies a common Gaussian majorant;
its range extends to $\sigma_M/M\asymp\sqrt{H_n}\to\infty$.
For $1/M\le|\theta|\le\theta_0$, the second estimate in
\eqref{eq:trunc-fourier-bound}, and $\tau_M\ge1$, give
$\left|\E e^{i\theta(\widetilde K-n)}\right|\le e^{-c n|\theta|^\alpha}$.
The normalized integral over this region is bounded by
\[
 C\sigma_M n^{-1/\alpha}
       \int_{n^{1/\alpha}/M}^\infty e^{-c v^\alpha}\dd v
 \le C H_n^{1/2-1/\alpha}
       \int_{H_n^{1/\alpha}}^\infty e^{-c v^\alpha}\dd v=o(1).
\]
This estimate is valid even when $H_n$ tends to infinity very slowly.
On $\theta_0\le|\theta|\le\pi$, the positive coefficient $a_1$ gives an
$e^{-cn}$ bound. Its integral, multiplied by $\sigma_M$, also tends to
zero, since the hypothesis bounds $M$ polynomially in $n$.
Fourier inversion now yields
$\sigma_M\Prob(\widetilde K=n)\to1/\sqrt{2\pi}$.
Inserting this in \eqref{eq:tilted-identity} proves \eqref{eq:hard}.
\end{proof}

\begin{remark}[Do not exponentiate an uncontrolled remainder]
Theorem~\ref{thm:hard} uses the \emph{exact} value $r_M$. It does not
permit the replacement
$r_M^{-n}=\rhoc^{-n}\exp(-nD_\alpha M^{-\alpha})(1+o(1))$
throughout its range. Theorem~\ref{thm:fold-drift} only controls
$\log(r_M/\rhoc)$ relatively; its unquantified error may be amplified by
$n$. That replacement needs a rate strong enough that
$n\left(\log(r_M/\rhoc)-D_\alpha M^{-\alpha}\right)\to0$.
The exact rearrangement \eqref{eq:radius-rearrangement} avoids this
problem in computation.
\end{remark}

\subsection{The two limits really differ}
For every fixed $M$, the power of $n$ in \eqref{eq:fixed-M} is $-3/2$,
and the exponential factor has radius $r_M>\rhoc$. For the untruncated
critical equation, the power is $-1-1/\alpha$ and the radius is $\rhoc$.
The drift $\log(r_M/\rhoc)\asymp M^{-\alpha}$ becomes order one after
multiplication by $n$ exactly at $M\asymp n^{1/\alpha}$. This agrees
with, but is not a replacement for, the full cutoff profile theorem:
the latter supplies both the missing amplitude and the conditioned
centering shift. The matching is therefore quantitative, not just a
comparison of formal critical exponents.

\section{Reproducible computation and numerical diagnostics}
\label{sec:computation}

\subsection{What was checked exactly}
The accompanying \texttt{code/verify.py} performs three independent
finite coefficient calculations: direct substitution into the nonlinear
equation, Lagrange's coefficient formula, and a sum over integer
partitions. For rational weights it compares the results through degree
$13$, using Python's exact \texttt{Fraction} arithmetic. The test weights
are $\beta/j^3$ with $\beta=1/3$, and their cutoff at $M=4$ with
$\beta=2/5$. The exponent $3$ is used to keep the arithmetic rational;
these tests verify the algebraic identity, which does not depend on the
restriction $1<\alpha<2$.

Additional exact comparisons check that retaining $M$ actions preserves
every coefficient through degree $M$. A symbolic calculation at
$\alpha=3/2$ verifies four coefficients of the Hahn implicit equation.
In total, \textbf{80 exact comparisons passed}. This is evidence for the
implementation and the displayed finite identities, not a substitute
for the asymptotic proofs.

For a single target sector $n$, the probability in
\eqref{eq:poisson} can be computed by a positive recurrence. If $p_k$ is
the probability generating coefficient at $t^k$, with the Poisson
parameter $n$ held fixed, then
\begin{equation}\label{eq:probability-recurrence}
 p_0=e^{-n\beta A_M(1)},\qquad
 p_k=\frac{n\beta}{k}\sum_{j=1}^{\min(k,M)}j a_jp_{k-j}.
\end{equation}
For the infinite family replace $A_M(1)$ by $A(1)$ and the upper summation
limit by $k$. This costs $O(n\min(n,M))$ arithmetic operations and $O(n)$
storage for one target $n$. The Poisson parameter changes when the target
$n$ changes; this complexity should not be mistaken for a simultaneous
all-sector algorithm of the same cost.

\subsection{The prototype with exponent \texorpdfstring{$\alpha=3/2$}{alpha=3/2}}
For $a=1$, $P=0$, and $\alpha=3/2$, the constants are
\begin{align*}
 \bc&=0.38279338399942656225\ldots,&
 \Bc&=0.90464481009458340321\ldots,\\
 \rhoc&=0.59839006994955853557\ldots,&
 g_{3/2}(0)&=0.24885478260493015222\ldots.
\end{align*}
They were evaluated with 75-digit arithmetic. The coefficient recurrence
was evaluated through target degree $4096$ using
\texttt{numpy.longdouble}; on platforms with a narrower exponent range,
the code switches to an arbitrary-precision recurrence if its initial
probability underflows. Fourier integrals were
computed in double precision, with two quadrature resolutions for the
cutoff density. None of these floating-point evaluations is an
outward-rounded interval certificate.

In Table~\ref{tab:critical}, $p_n=\Prob(K_n=n)$ and the corrected value is
the right side of \eqref{eq:threehalf-correction} without its remainder.
The disappearance of the nominal $n^{-1/3}$ term is visible in the
faster approach to the limiting constant.
\begin{table}[htbp]
\centering\small
\begin{tabular}{rrrr}
\toprule
$n$ & $b_n p_n$ & relative error & corrected value \\
\midrule
64 & 0.247067969 & -7.180e-03 & 0.247226460 \\
256 & 0.248168803 & -2.757e-03 & 0.248208582 \\
1024 & 0.248588377 & -1.071e-03 & 0.248598338 \\
4096 & 0.248750521 & -4.190e-04 & 0.248753012 \\
\bottomrule
\end{tabular}

\caption{Critical normalized coefficients. The relative error is measured
against $g_{3/2}(0)$; the correction is the first nonzero term allowed by
the cancellation sieve.}
\label{tab:critical}
\end{table}

Write $\Ra(L)=\Ra(L,0)$ at critical coupling. Table~\ref{tab:cutoff}
compares the limiting retained-coefficient fraction with the actual
ratio at $n=4096$ and $M=\lfloor Lb_n\rfloor$. The latter uses the rounded
integer cutoff, so $M/b_n$ is slightly less than the displayed target
$L$. At $L=1$ a substantial fraction is still lost. A larger fixed
multiple can be very accurate but, as Theorem~\ref{thm:budget} states,
has a strictly positive limiting loss.
\begin{table}[htbp]
\centering\small
\begin{tabular}{rrr}
\toprule
$L$ & $\mathcal R_{3/2}(L)$ & $n=4096$ ratio \\
\midrule
0.5 & 0.27103266 & 0.25750561 \\
1.0 & 0.76476600 & 0.76091153 \\
2.0 & 0.98459916 & 0.98537953 \\
4.0 & 0.99999962 & 0.99999988 \\
\bottomrule
\end{tabular}

\caption{Critical action-cutoff profile and finite-sector diagnostics.
The two quadrature resolutions differed by at most $2.34\times10^{-12}$
in the computed truncated densities; this difference is a diagnostic,
not a rigorous error bound.}
\label{tab:cutoff}
\end{table}

For the finite-fold constants, the calculations give
\begin{align*}
 v_{3/2}&=0.85403265659819698978\ldots,\\
 D_{3/2}&=0.59948018701656103250\ldots,\\
 V_{3/2}&=1.05291088528937168260\ldots.
\end{align*}
Table~\ref{tab:fold} illustrates the convergence in
Theorem~\ref{thm:fold-drift}. The complete data additionally contain the
scaled curvatures, critical-coupling-window comparisons at
$\lambda=-1,0,1$, and cutoff ratios at $n=256,1024,4096$.
\begin{table}[htbp]
\centering\small
\begin{tabular}{rrr}
\toprule
$M$ & $M\log\tau_M$ & $M^{3/2}\log(r_M/\rho_c)$ \\
\midrule
16 & 0.949469078 & 0.627376226 \\
64 & 0.906661654 & 0.619377140 \\
256 & 0.881417572 & 0.610854709 \\
1024 & 0.867971108 & 0.605513272 \\
\bottomrule
\end{tabular}

\caption{Moving finite-action folds. The two columns converge to
$v_{3/2}$ and $D_{3/2}$, respectively.}
\label{tab:fold}
\end{table}

\subsection{Reproduction and reliability boundaries}
Run
\begin{verbatim}
python -m pip install -r requirements.txt
python code/verify.py --max-n 4096
pdflatex -interaction=nonstopmode -halt-on-error article.tex
pdflatex -interaction=nonstopmode -halt-on-error article.tex
\end{verbatim}
A third LaTeX pass may be useful if page references move. The package
contains the generated CSV tables and \texttt{data/verification.json},
so typesetting does not require rerunning the computation. The formulas
and proofs are independent of the precision and quadrature choices in
the script. The numerical code does not verify global singularity
geometry, local-limit theorems, or the infinite support assertions by
finite sampling.

\section{Interpretation for transseries and a formalization route}
\label{sec:interpretation}

\subsection{Three scales, kept separate}
The model has three related but distinct asymptotic descriptions. First,
near $q=0$ the original $U_\beta(q)=\sum u_nq^n$ is an ordinary convergent
series; substituting $q=e^{-X}$ makes it a convergent integer-action
exponential transseries. Second, near its critical boundary one uses
$\varepsilon=\log(\rhoc/q)$, and $x$ has the fractional convergent Hahn
chart of Section~\ref{sec:hahn}. Third, as the original sector index
$n$ grows, the coefficients $u_n$ have the stable-density asymptotic
expansion of Section~\ref{sec:stable}. Theorem~\ref{thm:duality} relates
the last two descriptions. It does not identify their variables.

This separation resolves a potential ambiguity in statements about
``closure under inversion.'' The small-$q$ analytic germ is harmless;
the obstruction to a finite Puiseux chart occurs at a nonanalytic
boundary of the countable action sum. At rational $\alpha$ a finite
ramification exists and its minimal order is known exactly. At
irrational $\alpha$ the critical inverse still has a finite positive
Hahn grid, although it has no finite-sheeted Puiseux representation.
Thus finite-grid support and finite analytic ramification are different
properties.

\subsection{Finite-prefix universality has a precise boundary}
The finitely many altered weights change $d_0$, $d_1$, and the analytic
coefficients $d_k$. They therefore change the critical coupling, the
exponential growth factor, and the subleading sector corrections. They
do not change the leading critical exponents, the admissible Hahn grid,
the normalized cutoff profile, or the
arithmetic form of the cancellation sieve. The profile theorem is not
an assertion that all sequences with the same tail are identical;
it identifies exactly which quantities survive the normalization.

The positive-tail assumption supplies the Fourier damping. The condition
$a_1>0$ is a convenient sufficient aperiodicity hypothesis, not a claim
that more general lattice spans are impossible. Replacing it by a
nontrivial common divisor would require accessible residue classes and
corresponding Fourier aliases. Signed or complex weights would remove
the conditional probability interpretation altogether.

\subsection{What could be formalized first}
No new result in this article has been checked in Lean. The repository's
existing formalization of finite algebraic recurrences or residual
brackets does not, by itself, establish the analytic statements here.
A realistic staged interface is the following.

The first stage is finite algebra: the coefficient recurrence,
Lagrange's identity for a finite truncation, exact agreement through
degree $M$, and the finite multi-index support bounds. The second stage
is scalar special-function algebra: the gamma recurrence and reflection
identity proving \eqref{eq:duality-left}--\eqref{eq:duality-right}, with
reciprocal gamma used to include zeros cleanly. The third stage packages
the analytic implicit-function and Cauchy-majorant argument with explicit
inputs $H_0,m_\alpha,T_0,R_0$. These stages isolate checkable certificates
without pretending that symbolic coefficient equality proves analytic
convergence.

The probability stage should then formalize the independent Poisson
coefficient identity and the Fourier tail inequalities. Only after
those are in place should one assemble uniform local limits and
triangular-array limits. Their most useful reusable content is the
integrable-majorant interface, including uniform dependence on compact
parameter sets. Finally, the global analytic phase diagram can be linked
to a formal singularity-transfer theorem. This is a substantial program,
not an already completed formal crosswalk.

\section{Further research questions and concrete targets}
\label{sec:questions}

\paragraph{1. Slowly varying action tails.}
Replace the exact tail by $a_j\sim a j^{-1-\alpha}\ell(j)$ with slowly
varying $\ell$. Determine the correct implicit stable scale, the analogue
of the cutoff profile, and a usable inverse notation when de Bruijn
conjugates are needed. The finite positive Hahn grid may cease to be the
right representation; specify sufficient hypotheses that preserve it.
% ed. (2026-09-29): partial progress by a later sibling package (batch 48), at alpha = 2 only.
\begin{ednote}
Partial progress at the endpoint $\alpha=2$ only: the logarithmic-endpoint
package \cite{ed:lce} treats $a_j\sim a j^{-3}(\log j)^r$, $r\ge-1$, at
leading order (its \texttt{thm:weightedchart}, \texttt{thm:weightedbudget}).
% ed. (2026-09-29): on filing batch 50, "For 1<alpha<2 the question is open." was replaced: a later package answers it.
For $1<\alpha<2$ it is answered by the later package
\path{Slowly_Varying_Action_Tails_Critical_Transseries/} (batch~50; not
reviewed): scale, cutoff profile and budget for every slowly varying $\ell$
(its stable scale omits the factor $\Gamma(-\alpha)$ of $b_n$ here; with that
dictionary its profile at $\ell\equiv1$ is Theorem~\ref{thm:cutoff}); a
counterexample to the Hahn grid for $\ell=2+\sin(\log\log j)$; and a
convergent Lambert chart and an effective-index expansion for $(\log j)^m$.
\end{ednote}

\paragraph{2. Quantitative cutoff errors and certified action selection.}
Prove an explicit rate in Theorem~\ref{thm:cutoff}, uniform in a chosen
rectangle of $(L,\lambda)$, and enclose its density integrals by interval
arithmetic. The practical target is an algorithm which, given $n$ and a
relative tolerance, returns a certified cutoff $M$. The current profile
is exact as a limit, while the supplied quadrature is not certified.
% ed. (2026-09-29): partial progress by later sibling packages (batch 48), at alpha = 2 only.
\begin{ednote}
Partial progress at the endpoint $\alpha=2$ only. The logarithmic-endpoint
package \cite{ed:lce} proves an asymptotic error rate for its refined
cutoff law (its \texttt{thm:refined}) without explicit constants, and the
marginal package \cite{ed:mct} proves an explicit finite-$n$ lower bound on
the retained fraction (its \texttt{thm:finite-certificate}), whose recorded
evaluations use no directed rounding. Neither supplies the certified cutoff
algorithm asked for here; both list it as open. For $1<\alpha<2$ the
question is untouched.
\end{ednote}

\paragraph{3. An expanding coupling window.}
Find the largest growing range of $\lambda_n$ on which the all-order
expansion is relatively accurate, and match it to the fixed-subcritical
power-law regime and the supercritical saddle regime. This requires
large-argument estimates for stable-density derivatives and control of
how many correction terms are useful; the additive uniform theorem does
not already provide it.

\paragraph{4. Endpoint exponents and logarithmic crossover.}
Study joint limits $\alpha\downarrow1$ or $\alpha\uparrow2$ with $n$.
At the upper endpoint the diverging second-moment normalization competes
with Gaussian behavior; at the lower endpoint the mean itself is near
failure. Derive logarithmic critical charts and cutoff laws rather than
substituting an endpoint into formulas containing $\Gamma(-\alpha)$.
% ed. (2026-09-29): the upper-endpoint part of this question is answered by four later sibling packages (batches 48 and 49); the lower endpoint stays open.
\begin{ednote}
Answered at the upper endpoint by later packages of this series, which are
unrefereed research drafts; the lower endpoint $\alpha\downarrow1$ is treated
by none of them.
\emph{Fixed endpoint $\alpha=2$}, tail $a j^{-3}$ outside a finite prefix:
the logarithmic-endpoint package \cite{ed:lce} and the marginal package
\cite{ed:mct}, written independently of each other from a library copy of
this article, prove the same core. The pole of $\Gamma(-\alpha)$ cancels the
divergent quadratic coefficient and leaves $x^2\log(1/x)$; the critical
inverse is a convergent chart over a Lambert-$W_{-1}$ core; the sector
coefficients have an all-order Gaussian--logarithmic expansion; the retained
fraction tends to $e^{-1/(2s^2)}$ with the same first correction. Their
constants agree after a change of logarithmic normalization
($\ell_n=\log b_n+h$ in the first, $H_n=\log B_n$ in the second), and both
agree with the formal $\alpha\to2$ limits of this article's normal form,
budget (the omitted intensity $n\beta a/(\alpha M^\alpha)$ tends to
$1/(2s^2)$) and fold drift (Theorem~\ref{thm:fold-drift}).
\emph{Joint limit $\alpha_n\to2$}: answered, again independently, by the
confluent package \cite{ed:cct}, for $\varepsilon_n=2-\alpha_n\to0$ at any
rate and of either sign, at leading order, and by the stable--Gaussian
package \cite{ed:sge}, in the window $|2-\alpha_n|\log n\le K$, one order
further (a two-term uniform inverse, the first coefficient correction, and a
uniform first cutoff correction without finite-prefix moments). Where both
apply they agree term by term, and at $\alpha=2$ they reduce to the formulas
of the two fixed-endpoint packages.
\end{ednote}

\paragraph{5. Nonlinear slopes that break the exact substitution.}
For
$U=\beta\sum a_jq^j\exp((j+r_j)U)$ with $r_j=o(j)$, identify conditions
under which the stable boundary and budget persist. The parametrization
$t=qe^U$ is then no longer exact. A useful first case is
$r_j=\gamma j^\theta$, $0<\theta<1$, with a controlled complex domain.

\paragraph{6. Incommensurate actions and shrinking gaps.}
Replace integer actions $j$ by a genuinely incommensurate sequence, or
let a family of action gaps tend to zero. Determine whether the critical
inverse still admits a locally finite positive support semigroup and
which effective summability conditions replace the lattice Fourier
argument. This would address a different part of the repository's
parameter-dependent-action question.

\paragraph{7. Higher-order finite-fold drift.}
Expand $M\log\tau_M$, $\log(r_M/\rhoc)$, and $v_M$ beyond their leading
limits, allowing finite-prefix corrections. Establish a remainder strong
enough to exponentiate after multiplication by $n$ in a prescribed joint
range. This would turn the warning after Theorem~\ref{thm:hard} into a
usable sharp uniform formula.
% ed. (2026-09-29): answered at alpha = 2 only by later sibling packages (batch 48).
\begin{ednote}
Answered at the endpoint $\alpha=2$ only. The logarithmic-endpoint package
\cite{ed:lce} gives the finite-fold drift to all inverse-logarithmic orders
(its \texttt{thm:foldchart}) with an explicit condition for exponentiating it
(its \texttt{prop:exponentiation}); the marginal package \cite{ed:mct} gives
a convergent reduced chart for the drift (its \texttt{thm:finite-fold})
without an exponentiation criterion. For $1<\alpha<2$ the question is open.
\end{ednote}

\paragraph{8. Minimal support and deeper cancellations.}
The grid in \eqref{eq:hahn-grid} is sufficient, not always minimal.
Classify finite prefixes which suppress $d_2,d_3,\ldots$ and derive the
resulting first nonzero sector corrections. When rational $\alpha$
causes exponent collisions, seek an effective criterion for cancellation
of the entire grouped coefficient rather than only its gamma factor.

\paragraph{9. Divergent outer expansions and resurgence.}
Determine the large-order growth of the coefficients in
Theorem~\ref{thm:allorders}, optimal truncation, and possible Borel or
accelerated summation. Locate the relevant complex saddles and singular
values. Convergence of the local Hahn expansion and the coefficient
duality alone do not prove any of these claims.

\paragraph{10. Signed weights and competing complex singularities.}
Allow controlled negative or complex $a_j$ and replace positivity by
explicit continuation and saddle hypotheses. Determine when a
multi-saddle analogue of the cutoff profile exists. The nonnegative
Poisson conditioning proof will not survive unchanged, so an analytic
replacement is necessary.

\paragraph{11. Proof-producing computation.}
Implement a checker for finite support bounds, exact rational coefficient
identities, and local holomorphic majorants, followed by an interval
Fourier integrator. Connect its certificates to Lean statements with
separate hypotheses for analytic continuation, aperiodicity, and tail
bounds. The deliverable should distinguish a verified coefficient from a
verified asymptotic error estimate.

\section{Conclusion}
The decisive distinction is not merely between finitely many and
infinitely many formal symbols. It is whether the action sum remains
analytic at the boundary reached by the inverse. In the positive
power-tail class, the failure of a finite second moment changes the
critical exponent from $1/2$ to $1/\alpha$. Nevertheless, the inverse
still possesses a convergent, explicitly computable Hahn chart with
controlled support.

The same chart determines the all-order sector coefficients through an
exact gamma reflection identity. It explains why some apparently leading
corrections vanish, and why finite ramification is an arithmetic property
of $\alpha$. Truncating the countable action sum restores ordinary folds,
but preserving large sectors requires a sharp growing action budget.
The joint density profile supplies the missing transition, and the
finite-fold drift gives its complementary analytic explanation.

These results provide a rigorous model in which local inversion,
large-sector asymptotics, and action truncation can be compared without
interchanging limits informally. The broader questions of arbitrary
transseries realization, general resurgence, and noncommensurate action
families remain distinct research tasks.

\appendix
\section{Source audit and proof dependencies}
\label{app:sources}

\subsection{Repository snapshot and limits of inspection}
The repository sources used for the mathematical crosswalk were read at
commit
\begin{center}
\small\texttt{a795fcffa4b3ece22761d81fbed3c43be8b81795}.
\end{center}
The inspected group is
\begin{center}
\small\path{Analysis/Transseries/docs/series-and-transseries/}.
\end{center}
The group README's arrival descriptions and status caveats, the regularity
package README, and the moving-fold article's final research-question and
scope sections were inspected directly. The full canonical transseries
volume was not audited. In particular, the comparison above is not an
assertion that every possible overlap in that volume has been excluded.

During inspection, \texttt{main} advanced to
\begin{center}
\small\texttt{5804c7aff9955e6cbc09be2a71a295a07b238bc9},
\end{center}
whose parent is the pinned revision above. That later commit records a new incoming-archive
batch. This article does not claim to have reviewed those new archives.
The repository was not modified.

The group README describes the recent research packages as unmerged,
not claim-by-claim reviewed, and not Lean-formalized. Those qualifications
are preserved here rather than treating their statements as certified
premises. The underlying model and every theorem in this article are
proved directly.

\subsection{Dependency structure}
The analytic branch is
\begin{gather*}
 \text{polylogarithm expansion}\ \Longrightarrow\ \text{exact normal form}
 \ \Longrightarrow\\
 \text{convergent Hahn chart and coefficients}.
\end{gather*}
The sector branch is
\begin{gather*}
 \text{Lagrange identity}\ \Longrightarrow\ \text{Poisson coefficient representation}
 \ \Longrightarrow\\
 \text{Fourier local limits and corrections}.
\end{gather*}
Gamma reflection identifies the coefficients from these two branches.
The cutoff branch adds only tail Riemann sums, uniform Fourier bounds,
and the exact conditional identity. The finite-fold branch uses a
positive real root equation, its Riemann-sum limits, and a tilted Gaussian
local limit. Numerical verification is downstream of all four branches;
none uses numerical observations as a theorem hypothesis.

\section{A compact notation ledger}
\small
\begin{longtable}{@{}p{0.24\textwidth}p{0.70\textwidth}@{}}
\toprule
Symbol & Meaning\\
\midrule
\endhead
$\alpha,a,P$ & Tail exponent parameter, tail amplitude, and finite-prefix
polynomial in $A(t)=a\Li_{1+\alpha}(t)+P(t)$.\\[3pt]
$d_k$ & Analytic coefficients in \eqref{eq:d-def}; only $d_0,d_1$ are
ordinary finite moments.\\[3pt]
$G$ & $\Gamma(-\alpha)>0$.\\[3pt]
$\bc,\delta_\beta$ & Critical coupling $1/d_1$ and linear defect
$1-\beta d_1$.\\[3pt]
$\rho_\beta,B_\beta$ & Boundary value $e^{-\beta d_0}$ and fractional
coefficient $\beta aG$.\\[3pt]
$x,\varepsilon$ & $-\log t$ and $\log(\rho_\beta/q)$ near the boundary.\\[3pt]
$T,R,V$ & Hahn chart variables
$T=(\varepsilon/\Bc)^{1/\alpha}$, $R=T^{2-\alpha}$,
$x=TV(T,R)$.\\[3pt]
$\m,m_*,K,\Lambda$ & Multi-index, its total multiplicity, weighted index,
and Hahn exponent numerator in \eqref{eq:multiindex}.\\[3pt]
$b_n,\lambda_n$ & Stable sector scale and centered critical-window
coordinate in \eqref{eq:bn-lambda}.\\[3pt]
$g_\alpha,g_{\alpha,L}$ & Full and upper-jump-truncated centered stable
densities.\\[3pt]
$\nu_L,h_L$ & Omitted-jump intensity and omitted mean at cutoff $L$.\\[3pt]
$\Ra(L,\lambda)$ & Limiting retained-coefficient fraction in
\eqref{eq:cutoff-profile}.\\[3pt]
$\tau_M,r_M,v_M$ & Exact finite-action critical point, critical value,
and tilted curvature.\\[3pt]
$v_\alpha,D_\alpha,V_\alpha$ & Limiting finite-fold drift constants;
$v_\alpha$ is distinct from the curvature $v_M$.\\
\bottomrule
\end{longtable}
\normalsize

\begingroup
\small
\raggedright
\begin{thebibliography}{9}

\bibitem{dlmf}
NIST Digital Library of Mathematical Functions.
Section 25.12, \emph{Polylogarithms}, especially Eq.~25.12.12.
\url{https://dlmf.nist.gov/25.12}.
Consulted 29 September 2026.

\bibitem{flajolet}
P. Flajolet and R. Sedgewick.
\emph{Analytic Combinatorics}.
Cambridge University Press, 2009.
Author-maintained book site: \url{https://ac.cs.princeton.edu/}.

\bibitem{janson}
S. Janson.
Simply generated trees, conditioned Galton--Watson trees, random
allocations and condensation.
\emph{Probability Surveys} \textbf{9} (2012), 103--252.
\href{https://doi.org/10.1214/11-PS188}{doi:10.1214/11-PS188}.
Preprint \href{https://arxiv.org/abs/1112.0510}{arXiv:1112.0510};
see the critical infinite-variance discussion and Example~18.27.

\bibitem{chakrabarty}
A. Chakrabarty and G. Samorodnitsky.
Understanding heavy tails in a bounded world or, is a truncated heavy
tail heavy or not?
\emph{Stochastic Models} \textbf{28}(1) (2012), 109--143.
\href{https://doi.org/10.1080/15326349.2012.646551}{doi:10.1080/15326349.2012.646551}.
Preprint \href{https://arxiv.org/abs/1001.3218}{arXiv:1001.3218}.

\bibitem{repo-index}
VladimirReshetnikov/ProveIt repository.
\emph{Series and transseries}, group README and arrival-status descriptions.
Path \path{Analysis/Transseries/docs/series-and-transseries/README.md}.
Snapshot and inspected scope: Appendix~\ref{app:sources}.
\url{https://github.com/VladimirReshetnikov/ProveIt}.

\bibitem{repo-fold}
VladimirReshetnikov/ProveIt repository, research manuscript.
\emph{Critical Transseries at a Moving Fold: a uniform two-sheet atlas,
large-sector crossover, and resonant action splitting}.
Path below the preceding group:
\path{Critical_Transseries_Moving_Fold/article.tex}.
Final research-question and scope sections, source lines 1400--1580,
read at the pinned snapshot. September 2026.

\bibitem{repo-regularity}
VladimirReshetnikov/ProveIt repository, research package.
\emph{A Sharp Regularity Classification for Countable Exponential-Feedback
Transseries}.
Path below the group:
\path{Exponential_Feedback_Regularity_Classification/README.md}.
Package model, verification status, and limitations read at the pinned
snapshot. September 2026.

% ed. (2026-09-29): the four entries below are editorial additions for the notes on Questions 1, 2, 4 and 7.
\bibitem{ed:lce}
[Editorial addition, ProveIt, 2026-09-29.]
\emph{The Logarithmic Critical Endpoint: Convergent Lambert Charts,
All-Order Sector Laws, and Sharp Action-Cutoff Corrections}, research
package filed in batch~48. Path below the group:
\path{Logarithmic_Critical_Endpoint_Lambert_Charts/article.tex}.

\bibitem{ed:mct}
[Editorial addition, ProveIt, 2026-09-29.]
\emph{Marginal Critical Transseries: Convergent Lambert-W Charts, Universal
Cutoff Corrections, and Certified Action Budgets}, research package filed in
batch~48. Path below the group:
\path{Marginal_Critical_Transseries_Action_Budgets/Marginal_Critical_Transseries.tex}.

\bibitem{ed:cct}
[Editorial addition, ProveIt, 2026-09-29.]
\emph{Confluent Critical Transseries Across the Exponent-Two Boundary},
research package filed in batch~49. Path below the group:
\path{Confluent_Critical_Transseries_Exponent_Two_Boundary/Confluent_Critical_Transseries.tex}.

\bibitem{ed:sge}
[Editorial addition, ProveIt, 2026-09-29.]
\emph{Through the Stable--Gaussian Endpoint: Uniform Critical Transseries,
Prefix-Independent Cutoff Corrections, and Conditional Extremes}, research
package filed in batch~49. Path below the group:
\path{Stable_Gaussian_Endpoint_Uniform_Critical_Transseries/article.tex}.

\end{thebibliography}
\endgroup
\end{document}
